\problemname{}

%\illustration{0.3}{image.jpg}{
%    Caption of the illustration (optional).
%    CC BY-SA 4.0 by X on \href{https://example.com/reference-to-image}{Y}
%}

The process of forming the Dutch government has taken more than half a year for three elections in a row. Perhaps we can streamline the initial stages of coalition building?

The first step after the election results is to find a group of parties (called a \emph{coalition}) with enough seats to have a strict majority.
Your task is to count the number of candidate coalitions that satisfy specific conditions.
A coalition is considered a \emph{candidate coalition} if it meets these two criteria:

\begin{description}
    \item[Strict Majority:] The total number of seats held by the coalition must be strictly more than half of the total seats across all parties.
    \item[No Superfluous Parties:] The coalition must be minimal in the sense that removing any one party from the coalition would cause it to lose its strict majority.
\end{description}

\begin{figure}[!h]
    \centering
    \includegraphics[width=0.5\textwidth]{sample2.pdf}
    \caption{Illustration of Sample Input 2.}
    \label{fig:sample2}
\end{figure}

\begin{Input}
    The input consists of:
    \begin{itemize}
        \item One line with an integer $n$ ($1 \le n \le 60$), the number of parties.
        \item One line with $n$ integers $p$ ($1 \le p \le 10\,000$), the number of seats each party has.
    \end{itemize}
\end{Input}

\begin{Output}
     Output the total number of candidate coalitions that satisfy the criteria above.
\end{Output}
